\chapter{Introducción}

La gestión de residuos es un desafío crítico en el contexto del cambio climático y la sostenibilidad ambiental debido a su impacto directo en la generación de gases de efecto invernadero (GEI), como el metano (CH$4$) y el dióxido de carbono (CO$2$) \cite{filonchyk2024greenhouse}
\parencite{filonchyk2024greenhouse} La gestión de residuos es un desafío crítico en el contexto del cambio climático y la sostenibilidad ambiental debido a su impacto directo en la generación de gases de efecto invernadero (GEI), como el metano (CH$4$) y el dióxido de carbono (CO$2$) (Filonchyk, et al., 2024; Kabange, et al., 2023).   Estos residuos también contribuyen a la contaminación del suelo y el agua, afectando negativamente la biodiversidad y la salud humana (FAO, 2021; IPCC, 2014; Leddin, 2024).  El sector agrícola, responsable de aproximadamente el 24 \% de las emisiones globales de GEI, enfrenta este reto principalmente por prácticas inadecuadas en la gestión de residuos, fermentación entérica y uso excesivo de fertilizantes sintéticos (Smith et al., 2014). Esta situación demanda la implementación de estrategias sostenibles e innovadoras para la correcta valorización de residuos.
